\section{Hemalurgic Constructs}
\label{sec:hemalurgyConstructs}

Using the forbidden art, a skillful Hemalurgist may not only steal powers from their victims, but also use the spikes to create twisted monstrosities. Example might include koloss warriors or hemalurgic chimeras. The process requires much time and skill, but a player may work with the GM to determine how to approach it.

\subsubsection{Time-consuming Process}

It takes a lot of time and effort to properly study this branch of Hemalurgy, prepare all required spikes, spike the victim and wait for the effect to fully take place. This is not something that can be quickly completed in a midst of Combat and instead should be treated as a goal the PC pursues.

\subsubsection{Limited to a Reward}

An ability to create a limitless army of hemalurgic monsters would quickly spiral out of control. To limit it, the creation of such constructs can be only achieved as a result of a reward. Refer to "Reward Categories - Companions" (\handbookMistborn{}, page 286.) for more details.

\subsubsection{Before Starting...}

At the moment a PC starts a goal to create a Hemalurgic Construct, they must define the desired result they are trying to achieve. They can choose one of the better-known monstrosities, or use their creativity and design something on their own.

When designing a new Construct, the player may wish to apply some of the below effects or to define some other, more unique powers they wish to bestow upon their creation:

\begin{itemize}
	\item Increased Strength, Speed, or other attributes
	\item Increased Deflect, Focus or Investiture
	\item Access to Invested Arts
	\item Actions or Features of a wild animal
	\item Special abilities
\end{itemize}

Depending on the desired results, the PC will need to gather required Hemalurgic Spikes to complete their creation. Work with your GM to clearly define the requirements. Note that for more complex Constructs, you might need unique metals, like trellium or atium.

\subsubsection{Known Constructs}

Here is the list of official Hemalurgic Constructs presented in \cosmereRPG{} sources. You can pick one of them as a final result of your creation or use them as an inspiration for more personalized experiments:

\begin{itemize}
	\item \textbf{Hemalurgic Beast} \\ Tier 1 Boss \\ see \bookMistbornLegacy{}, page 218.
	\item \textbf{Koloss Neophyte} \\ Tier 2 Riva \\ see \bookWorldGuideMistborn{}, page 244.
	\item \textbf{Hemalurgic Chimera} \\ Tier 3 Minion \\ see \bookWorldGuideMistborn{}, page 234.	
\end{itemize}


\subsubsection{Stats' Range}

When creating your own Hemalurgic Construct make sure to work with your GM to confirm the final design is well balanced. Compare with other \cosmereRPG{} adversaries for reference or follow the guidelines below:

\begin{itemize}
	\item \textbf{Tier 1 Minion} \\ Max attribute: 3 \\ Total attributes: 6-7 \\ Max skill rank: 3 \\ Total skill ranks: 7-8 \\ Health: 10-13
	
	\item \textbf{Tier 1 Rival / Tier 2 Minion} \\ Max attribute: 4 (tier 1) - 5 (tier 2) \\ Total attributes: 14-15 \\ Max skill rank: 3 \\ Total skill ranks: 12-15 \\ Health: 20-25
	
	\item \textbf{Tier 2 Rival / Tier 3 Minion} \\ Max attribute: 6 \\ Total attributes: 17-20 \\ Max skill rank: 4 \\ Total skill ranks: 15-20 \\ Health: 50-60 
\end{itemize}


\subsubsection{Construct's Progression}

As your PC advances to a higher tier, your Construct may advance as well as a reward. For example, your \textbf{Koloss Neophyte} may grow to become \textbf{Koloss Veteran}.

